\documentclass[11pt,a4paper]{article}

% Packages
\usepackage[utf8]{inputenc}
\usepackage[T1]{fontenc}
\usepackage{CJKutf8}
\usepackage{geometry}
\geometry{margin=2.5cm}
\usepackage{graphicx}
\usepackage{booktabs}
\usepackage{hyperref}
\usepackage{amsmath}
\usepackage{caption}
\usepackage{subcaption}
\usepackage{float}
\usepackage{enumitem}
\usepackage{xcolor}
\usepackage{listings}

\hypersetup{
    colorlinks=true,
    linkcolor=blue,
    urlcolor=blue,
    citecolor=blue
}

\lstset{
    basicstyle=\ttfamily\small,
    breaklines=true,
    frame=single,
    backgroundcolor=\color{gray!10}
}

\title{\textbf{HW2: Detect AI Generated Text}\\[0.3em]
\large RNN and Transformer — Comprehensive Report}
\author{\begin{CJK*}{UTF8}{gbsn}李捷新\end{CJK*}}
\date{April 2026}

\begin{document}

\maketitle

\begin{abstract}
This report presents a comprehensive study on detecting AI-generated text using BERT-based classifiers. We fine-tune both \texttt{bert-base-cased} (108M parameters) and \texttt{bert-large-cased} (334M parameters) on the DAIGT V2 Train Dataset (44,868 essays). BERT-large achieves perfect ROC-AUC of 1.0000 and accuracy of 0.9980. We compare against TF-IDF baselines (best AUC: 0.9997 with Linear SVM), conduct extensive ablation studies on learning rate, batch size, sequence length, and training epochs, and perform adversarial attacks using a locally deployed DeepSeek-R1:8B model via Ollama plus programmatic perturbations. The BERT detector proves highly robust: 0/30 single-pass LLM attacks succeed, 1/5 iterative attacks succeed (20\%), and only 1/30 programmatic perturbations evade detection (3.3\%). All experiments were conducted on an NVIDIA RTX 4090 GPU with fp16 mixed precision.
\end{abstract}

\tableofcontents
\newpage

% ============================================================
\section{Data Analysis \& Baseline (Part 1)}
% ============================================================

\subsection{Exploratory Data Analysis}

We analyzed four text-level features comparing human-written (Label 0) and AI-generated (Label 1) essays. Statistical significance was evaluated using the Mann-Whitney U test, with effect sizes measured by Cohen's $d$.

Here, Cohen's $d$ is computed as $(\mu_{\text{human}} - \mu_{\text{AI}}) / \sigma_{\text{pooled}}$, so positive values mean the feature is larger for human essays, while negative values mean it is larger for AI essays.

\begin{table}[H]
\centering
\caption{EDA: Feature comparison between human and AI text}
\label{tab:eda}
\begin{tabular}{lcccc}
\toprule
\textbf{Feature} & \textbf{Human (mean)} & \textbf{AI (mean)} & \textbf{Cohen's $d$} & \textbf{$p$-value} \\
\midrule
Avg Word Length     & 4.53  & 5.09  & $-0.733$ & $\approx 0$ \\
Word Count          & 418.3 & 329.4 & $+0.594$ & $\approx 0$ \\
Sentence Count      & 22.3  & 19.1  & $+0.384$ & $7.11 \times 10^{-217}$ \\
Type-Token Ratio    & 0.4785 & 0.5104 & $-0.366$ & $1.99 \times 10^{-229}$ \\
\bottomrule
\end{tabular}
\end{table}

All four features show statistically significant differences ($p \approx 0$). The strongest signal is \textbf{average word length} (Cohen's $d = -0.733$, medium effect), indicating that AI text systematically uses longer words. Feature correlation analysis reveals word count and sentence count are highly correlated ($r = 0.80$), while TTR is negatively correlated with word count ($r = -0.59$).

\begin{figure}[H]
\centering
\includegraphics[width=0.85\textwidth]{results_baseline/eda_distributions.png}
\caption{Distribution comparison of text features between human and AI essays.}
\label{fig:eda}
\end{figure}

\begin{figure}[H]
\centering
\includegraphics[width=0.85\textwidth]{results_baseline/enhanced_eda.png}
\caption{Feature correlation heatmap and all-baselines ROC comparison.}
\label{fig:enhanced_eda}
\end{figure}

\subsection{Baseline Models}

We implemented four classifiers using TF-IDF vectorization (max\_features=10{,}000, ngram\_range=(1,2)):

\begin{table}[H]
\centering
\caption{Baseline model performance (TF-IDF features)}
\label{tab:baseline}
\begin{tabular}{lccccc}
\toprule
\textbf{Model} & \textbf{ROC-AUC} & \textbf{Accuracy} & \textbf{Precision} & \textbf{Recall} & \textbf{F1} \\
\midrule
Logistic Regression      & 0.9993 & 0.9939 & 0.9980 & 0.9863 & 0.9921 \\
\textbf{Linear SVM}      & \textbf{0.9997} & \textbf{0.9970} & 0.9971 & 0.9951 & 0.9961 \\
Multinomial Naive Bayes  & 0.9955 & 0.9692 & 0.9781 & 0.9423 & 0.9598 \\
Random Forest (200 trees)& 0.9987 & 0.9884 & 0.9968 & 0.9734 & 0.9850 \\
\bottomrule
\end{tabular}
\end{table}

Linear SVM achieves the highest baseline AUC (0.9997), approaching BERT-level performance. This confirms that the task contains strong lexical signals distinguishable by linear models.

\begin{figure}[H]
\centering
\includegraphics[width=0.85\textwidth]{results_baseline/baseline_confusion_matrices.png}
\caption{Confusion matrices for all four baseline classifiers.}
\label{fig:baseline_cm}
\end{figure}

% ============================================================
\section{BERT Fine-Tuning \& Scaling (Part 2)}
% ============================================================

\subsection{Experimental Setup}

\begin{itemize}[nosep]
\item \textbf{Models}: \texttt{bert-base-cased} (108M params), \texttt{bert-large-cased} (334M params)
\item \textbf{Tokenizer}: \texttt{AutoTokenizer.from\_pretrained()} with matching checkpoint
\item \textbf{Training}: HuggingFace \texttt{Trainer} API, fp16 mixed precision, AdamW optimizer
\item \textbf{Hyperparameters}: learning\_rate = $2 \times 10^{-5}$, weight\_decay = 0.01, warmup\_ratio = 0.1
\item \textbf{Data split}: 80/20 stratified train/validation
\item \textbf{Hardware}: NVIDIA RTX 4090 (24GB VRAM)
\end{itemize}

\subsection{Core Experiments}

\begin{table}[H]
\centering
\caption{Core BERT fine-tuning results}
\label{tab:bert_core}
\begin{tabular}{lccccr}
\toprule
\textbf{Experiment} & \textbf{Params} & \textbf{ROC-AUC} & \textbf{Accuracy} & \textbf{Batch} & \textbf{Time} \\
\midrule
BERT-base (3ep, len=512)  & 108M & 0.9999 & 0.9947 & 32 & 8.6 min \\
BERT-base (3ep, len=256)  & 108M & 0.9998 & 0.9903 & 64 & 4.0 min \\
BERT-base (5ep, len=512)  & 108M & 0.9999 & 0.9957 & 32 & 14.4 min \\
BERT-large (3ep, len=512) & 334M & \textbf{1.0000} & \textbf{0.9980} & 16 & 27.2 min \\
\bottomrule
\end{tabular}
\end{table}

\textbf{BERT-large achieves perfect ROC-AUC (1.0000)} and the highest accuracy (0.9980) across all experiments.

\subsection{Learning Rate Ablation}

We train BERT-base with three learning rates while holding other hyperparameters fixed (3 epochs, max\_len=512, batch\_size=32):

\begin{table}[H]
\centering
\caption{Learning rate ablation (BERT-base, 3 epochs)}
\label{tab:lr_ablation}
\begin{tabular}{cccc}
\toprule
\textbf{Learning Rate} & \textbf{ROC-AUC} & \textbf{Accuracy} & \textbf{Time} \\
\midrule
$1 \times 10^{-5}$ & 0.9998 & 0.9961 & 8.6 min \\
$\mathbf{2 \times 10^{-5}}$ & \textbf{0.9999} & 0.9947 & 8.6 min \\
$5 \times 10^{-5}$ & 0.9999 & 0.9924 & 8.6 min \\
\bottomrule
\end{tabular}
\end{table}

\textbf{Finding:} $\text{lr}=2\times10^{-5}$ achieves the best AUC. Lower learning rate ($1\times10^{-5}$) converges slower, yielding slightly lower AUC despite higher accuracy. Higher learning rate ($5\times10^{-5}$) leads to lower accuracy due to overshoot.

\begin{figure}[H]
\centering
\includegraphics[width=0.75\textwidth]{results_bert-base-cased/lr_ablation.png}
\caption{Learning rate ablation: ROC-AUC and Accuracy comparison.}
\label{fig:lr_ablation}
\end{figure}

\subsection{Batch Size Ablation}

We train BERT-base with three batch sizes (3 epochs, max\_len=512, $\text{lr}=2\times10^{-5}$):

\begin{table}[H]
\centering
\caption{Batch size ablation (BERT-base, 3 epochs)}
\label{tab:bs_ablation}
\begin{tabular}{cccc}
\toprule
\textbf{Batch Size} & \textbf{ROC-AUC} & \textbf{Accuracy} & \textbf{Time} \\
\midrule
8   & 0.9999 & \textbf{0.9974} & 11.4 min \\
16  & 0.9999 & 0.9954 & 9.1 min \\
32  & 0.9999 & 0.9947 & 8.6 min \\
\bottomrule
\end{tabular}
\end{table}

\textbf{Finding:} All batch sizes achieve the same ROC-AUC (0.9999). Smaller batch size (BS=8) yields the highest accuracy (0.9974) due to implicit regularization from noisier gradients, at the cost of 33\% longer training time.

\begin{figure}[H]
\centering
\includegraphics[width=0.75\textwidth]{results_bert-base-cased/batch_size_ablation.png}
\caption{Batch size ablation: ROC-AUC and Accuracy comparison.}
\label{fig:bs_ablation}
\end{figure}

\subsection{BERT-large Epoch Analysis}

\begin{table}[H]
\centering
\caption{BERT-large: 3 vs 5 epochs}
\label{tab:large_epoch}
\begin{tabular}{cccc}
\toprule
\textbf{Epochs} & \textbf{ROC-AUC} & \textbf{Accuracy} & \textbf{Time} \\
\midrule
3 & \textbf{1.0000} & \textbf{0.9980} & 27.2 min \\
5 & 0.9998 & 0.9960 & 45.3 min \\
\bottomrule
\end{tabular}
\end{table}

\textbf{Finding:} BERT-large with 5 epochs \textbf{degrades} compared to 3 epochs (AUC: $1.0000 \to 0.9998$, Acc: $0.9980 \to 0.9960$). The larger model converges quickly and additional training causes overfitting.

\subsection{Comprehensive Comparison (All 9 Experiments)}

\begin{table}[H]
\centering
\caption{All nine BERT experiments}
\label{tab:all_bert}
\begin{tabular}{clcccr}
\toprule
\textbf{\#} & \textbf{Configuration} & \textbf{ROC-AUC} & \textbf{Accuracy} & \textbf{BS} & \textbf{Time} \\
\midrule
1 & base, 3ep, 512, lr=2e-5    & 0.9999 & 0.9947 & 32 & 8.6m \\
2 & large, 3ep, 512, lr=2e-5   & \textbf{1.0000} & \textbf{0.9980} & 16 & 27.2m \\
3 & base, 3ep, 256, lr=2e-5    & 0.9998 & 0.9903 & 64 & 4.0m \\
4 & base, 5ep, 512, lr=2e-5    & 0.9999 & 0.9957 & 32 & 14.4m \\
5 & base, 3ep, 512, lr=1e-5    & 0.9998 & 0.9961 & 32 & 8.6m \\
6 & base, 3ep, 512, lr=5e-5    & 0.9999 & 0.9924 & 32 & 8.6m \\
7 & base, 3ep, 512, bs=8       & 0.9999 & 0.9974 & 8  & 11.4m \\
8 & base, 3ep, 512, bs=16      & 0.9999 & 0.9954 & 16 & 9.1m \\
9 & large, 5ep, 512, lr=2e-5   & 0.9998 & 0.9960 & 16 & 45.3m \\
\bottomrule
\end{tabular}
\end{table}

\begin{figure}[H]
\centering
\includegraphics[width=0.95\textwidth]{results_bert-base-cased/all_experiments_comparison.png}
\caption{All 9 BERT experiments: ROC-AUC and Accuracy comparison.}
\label{fig:all_experiments}
\end{figure}

\subsection{Per-Class Error Analysis}

Using the best model (BERT-large, 3 epochs):

\begin{table}[H]
\centering
\caption{Confusion matrix for BERT-large (3ep)}
\label{tab:error}
\begin{tabular}{lcccc}
\toprule
& \textbf{TP} & \textbf{TN} & \textbf{FP} & \textbf{FN} \\
\midrule
Count & 3,492 & 5,464 & 10 & 8 \\
Rate  & ---   & ---   & 0.18\% & 0.23\% \\
\bottomrule
\end{tabular}
\end{table}

Only 18 total misclassifications out of 8,974 validation samples.

\subsection{Text Length Impact}

\begin{table}[H]
\centering
\caption{Accuracy by essay length quartile (BERT-large)}
\label{tab:text_length}
\begin{tabular}{lccc}
\toprule
\textbf{Quartile} & \textbf{$N$} & \textbf{Accuracy} & \textbf{Errors} \\
\midrule
Q1 (short, $\leq$277 words)  & 2,248 & 0.9969 & 7 \\
Q2 (278--354 words)           & 2,250 & 0.9982 & 4 \\
Q3 (355--453 words)           & 2,237 & 0.9973 & 6 \\
Q4 (long, $>$453 words)       & 2,239 & \textbf{0.9996} & 1 \\
\bottomrule
\end{tabular}
\end{table}

\textbf{Finding:} Longer essays are classified more accurately (Q4: 0.9996 vs Q1: 0.9969). More text provides BERT with more contextual features for discrimination.

\begin{figure}[H]
\centering
\includegraphics[width=0.85\textwidth]{results_bert-large-cased/error_analysis.png}
\caption{Per-class error analysis and text length impact.}
\label{fig:error_analysis}
\end{figure}

\subsection{Scaling Analysis}

\begin{table}[H]
\centering
\caption{Scaling: BERT-base vs BERT-large}
\label{tab:scaling}
\begin{tabular}{lcccc}
\toprule
\textbf{Model} & \textbf{Parameters} & \textbf{ROC-AUC} & \textbf{Accuracy} & \textbf{Time} \\
\midrule
BERT-base  & 108M & 0.9999 & 0.9947 & 8.6 min \\
BERT-large & 334M & 1.0000 & 0.9980 & 27.2 min \\
\midrule
\textbf{Improvement} & $+3.1\times$ & $+0.0001$ & $+0.33$ pp & $+3.2\times$ \\
\bottomrule
\end{tabular}
\end{table}

BERT-large achieves a marginal improvement (AUC $+0.0001$, Accuracy $+0.33$ percentage points) at $3.2\times$ the training cost. The task exhibits \textbf{diminishing returns} from scaling --- even BERT-base nearly saturates. However, BERT-large achieves perfect AUC (1.0000) and is preferred when compute budget allows.

\begin{figure}[H]
\centering
\includegraphics[width=0.85\textwidth]{results_bert-base-cased/comparison_curves.png}
\caption{Training and evaluation loss curves for BERT-base and BERT-large.}
\label{fig:comparison}
\end{figure}

\begin{figure}[H]
\centering
\includegraphics[width=0.85\textwidth]{results_bert-base-cased/bert_confusion_confidence.png}
\caption{Confusion matrices and confidence distributions for BERT-base and BERT-large.}
\label{fig:bert_cm}
\end{figure}

\begin{figure}[H]
\centering
\includegraphics[width=0.7\textwidth]{results_bert-base-cased/epoch_ablation.png}
\caption{Epoch ablation: Training loss comparison for 3-epoch vs 5-epoch BERT-base.}
\label{fig:epoch_ablation}
\end{figure}

% ============================================================
\section{Adversarial Attack with Local LLM (Part 3)}
% ============================================================

\subsection{Setup}

\begin{itemize}[nosep]
\item \textbf{Local LLM}: DeepSeek-R1:8B (\texttt{deepseek-r1:8b-llama-distill-fp16}) deployed via Ollama
\item \textbf{Detector}: BERT-large-cased (3 epochs, max\_len=512), AUC = 1.0000
\item \textbf{Essays}: 10 human-written essays from the validation set
\item \textbf{Strategies}: Academic rewriting, Stylistic rewriting, Paraphrase
\end{itemize}

\subsection{LLM Single-Pass Attack Results}

\begin{table}[H]
\centering
\caption{Single-pass adversarial attack results (10 essays $\times$ 3 strategies = 30 attacks)}
\label{tab:adversarial}
\begin{tabular}{lccc}
\toprule
\textbf{Strategy} & \textbf{Attacks} & \textbf{Fooled} & \textbf{Avg $P(\text{AI})$} \\
\midrule
Academic    & 10 & 0 & 0.9996 \\
Stylistic   & 10 & 0 & 0.9993 \\
Paraphrase  & 10 & 0 & 0.9997 \\
\midrule
\textbf{Total} & \textbf{30} & \textbf{0 (0\%)} & \textbf{0.9998} \\
\bottomrule
\end{tabular}
\end{table}

\textbf{Result:} All 30 single-pass LLM rewrites were caught. The detector assigns AI probability $\approx 1.0$ to all rewrites with very high confidence.

\subsection{LLM Iterative Attack Results}

We implemented a feedback-loop attack: after each rewrite, we report the AI probability to the LLM and ask it to further revise. Results for 5 essays with up to 3 iterations:

\begin{table}[H]
\centering
\caption{Iterative adversarial attack results (3-round feedback loop)}
\label{tab:iterative}
\begin{tabular}{lccccl}
\toprule
\textbf{Essay} & \textbf{Original} & \textbf{Round 1} & \textbf{Round 2} & \textbf{Round 3} & \textbf{Result} \\
\midrule
1 & 0.0000 & 1.0000 & 1.0000 & 1.0000 & CAUGHT \\
2 & 0.0000 & 1.0000 & 1.0000 & 1.0000 & CAUGHT \\
3 & 0.0000 & 1.0000 & 1.0000 & 1.0000 & CAUGHT \\
4 & 0.0000 & 0.9746 & 1.0000 & 1.0000 & CAUGHT \\
5 & 0.0000 & \textbf{0.1250} & --- & --- & \textbf{FOOLED} \\
\bottomrule
\end{tabular}
\end{table}

\textbf{Result:} 1/5 (20\%) fooled. Essay 5 dropped to $P(\text{AI})=0.125$ after one iteration, successfully evading the detector.

\subsection{Adversarial Text Example}

Below is a concrete comparison for Essay~5 (index 5258), the only successfully fooled essay in the iterative attack.

\subsubsection*{Original Human-Written Text ($P(\text{AI})=0.0003$)}

\begin{quote}
``What would you think if you heard `Your principal has decided that all students must participate in at least one extracurricular activity?' 
I agree with this prompt because kids wont get in trouble as much. Also, Some kids dont have many friends. It can also be a good life-trait.

Firstly, Kids make it look like they love being in trouble. If we had an after school activity they liked to do, They would have a reason to stay out of trouble. Generic\_Name, was always in trouble in school. Until, The\ldots''
\end{quote}

\textbf{Analysis of original:} Characteristic human-writing markers are present: grammatical errors (``wont'', ``dont''), informal punctuation, inconsistent capitalisation, and first-person anecdote. BERT assigns $P(\text{AI})=0.0003$ with very high confidence.

\subsubsection*{LLM-Rewritten Text (1 Iteration, $P(\text{AI})=0.125$) — FOOLED}

\begin{quote}
``The idea that every student must participate in at least one extracurricular activity seems like a solid plan to me. I think it's a smart move by the principal because it addresses a few key issues that many schools face today. Students who are involved in activities outside the classroom tend to stay out of trouble, build social connections, and develop skills that prepare them for life\ldots''
\end{quote}

\textbf{Why this fooled BERT:} The rewrite preserves the first-person voice and retains the student-like perspective (``seems like a solid plan to me''), while correcting grammar errors. Crucially, the LLM did \emph{not} introduce formal academic structure or topic-sentence paragraphs. The result reads like a thoughtful student essay rather than a polished AI composition, causing BERT's confidence to drop from $P(\text{AI})\approx1.0$ (for typical LLM rewrites) to $P(\text{AI})=0.125$.

\subsubsection*{Comparison with a Failed Attack (Essay~1)}

\begin{quote}
\textbf{Original:} ``You know, I've been thinking about this whole phone thing at school. It's like, on one hand, phones are super cool\ldots''
\end{quote}

After LLM rewrite: $P(\text{AI})=0.9999$ (CAUGHT). The LLM introduced a bold heading (``\textbf{Should Phones Be Allowed in School?}'') and structured paragraphs — both strong AI-text signals that caused the detector to become \emph{more} confident.

\textbf{Key Insight} (the rewriting paradox): For most essays, LLM rewriting \emph{replaces one AI signature with another}, driving $P(\text{AI})\to1.0$. Essay~5 succeeded precisely because the model happened to produce a colloquial, error-corrected but conversational rewrite that preserved human-like discourse patterns.

\subsection{Programmatic Adversarial Attacks (No LLM)}

We test simple text perturbations on 10 AI-generated essays to evaluate detector robustness against surface-level modifications:

\begin{table}[H]
\centering
\caption{Programmatic adversarial attack results (10 AI essays $\times$ 3 strategies)}
\label{tab:programmatic}
\begin{tabular}{lcccc}
\toprule
\textbf{Strategy} & \textbf{Attacks} & \textbf{Evaded} & \textbf{Rate} & \textbf{Avg $P(\text{AI})$ after} \\
\midrule
Typo Injection (2\% char swap)    & 10 & 1 & 10\% & 0.9159 \\
Homoglyph Substitution (3\%)     & 10 & 0 & 0\%  & 0.9999 \\
Whitespace Perturbation (5\%)    & 10 & 0 & 0\%  & 1.0000 \\
\midrule
\textbf{Total} & \textbf{30} & \textbf{1 (3.3\%)} & --- & --- \\
\bottomrule
\end{tabular}
\end{table}

\textbf{Finding:} Simple text perturbations are largely ineffective. Only typo injection had a minor effect (1/10 evasion), reducing average AI probability to 0.9159. Homoglyph and whitespace attacks had virtually no impact, showing BERT's tokenizer is robust to these surface-level perturbations.

\subsection{All Attack Methods Comparison}

\begin{table}[H]
\centering
\caption{Summary of all adversarial attack methods}
\label{tab:attack_summary}
\begin{tabular}{lccc}
\toprule
\textbf{Method} & \textbf{Attacks} & \textbf{Fooled/Evaded} & \textbf{Success Rate} \\
\midrule
LLM Single-Pass     & 30 & 0 & 0.0\% \\
Programmatic         & 30 & 1 & 3.3\% \\
LLM Iterative        & 5 ($\times$3 rounds) & 1 & 20.0\% \\
\bottomrule
\end{tabular}
\end{table}

\subsection{Why BERT is Robust}

\begin{enumerate}
\item \textbf{Deep semantic understanding}: BERT captures paragraph flow, discourse structure, and vocabulary distribution patterns that survive single-pass rewrites.
\item \textbf{LLM-rewriting paradox}: Using an LLM to rewrite text \emph{replaces one AI signature with another}. Human essays become \emph{more} AI-like after rewriting ($P(\text{AI})$: $\sim\!0.0000 \to \sim\!1.0000$).
\item \textbf{Training diversity}: The DAIGT V2 dataset contains AI text from multiple LLMs, making the detector generalize well across generator styles.
\item \textbf{Tokenizer robustness}: Homoglyph and whitespace perturbations are neutralized by BERT's WordPiece tokenizer, which maps perturbed characters to \texttt{[UNK]} tokens or normalizes whitespace.
\end{enumerate}

\begin{figure}[H]
\centering
\includegraphics[width=0.85\textwidth]{results_adversarial/adversarial_attack_summary.png}
\caption{Adversarial attack summary: AI probability shift and fool rate by strategy.}
\label{fig:adversarial}
\end{figure}

\begin{figure}[H]
\centering
\includegraphics[width=0.85\textwidth]{results_adversarial/enhanced_adversarial_analysis.png}
\caption{Iterative attack AI probability trajectory and analysis.}
\label{fig:adversarial_analysis}
\end{figure}

\begin{figure}[H]
\centering
\includegraphics[width=0.85\textwidth]{results_adversarial/programmatic_attacks.png}
\caption{Programmatic adversarial attacks: AI probability before and after perturbation.}
\label{fig:programmatic}
\end{figure}

% ============================================================
\section{Summary \& Conclusions}
% ============================================================

\begin{table}[H]
\centering
\caption{Final performance comparison across all models}
\label{tab:summary}
\begin{tabular}{llcc}
\toprule
\textbf{Model} & \textbf{Type} & \textbf{ROC-AUC} & \textbf{Accuracy} \\
\midrule
TF-IDF + Logistic Regression & Baseline        & 0.9993 & 0.9939 \\
TF-IDF + Linear SVM          & Baseline        & 0.9997 & 0.9970 \\
TF-IDF + Random Forest        & Baseline        & 0.9987 & 0.9884 \\
TF-IDF + Naive Bayes          & Baseline        & 0.9955 & 0.9692 \\
BERT-base (3ep, bs=32)        & Deep Learning   & 0.9999 & 0.9947 \\
BERT-base (3ep, bs=8)         & Deep Learning   & 0.9999 & 0.9974 \\
BERT-base (5ep, bs=32)        & Deep Learning   & 0.9999 & 0.9957 \\
\textbf{BERT-large (3ep)}     & Deep Learning   & \textbf{1.0000} & \textbf{0.9980} \\
BERT-large (5ep)              & Deep Learning   & 0.9998 & 0.9960 \\
\bottomrule
\end{tabular}
\end{table}

\textbf{Key findings}:
\begin{enumerate}
\item \textbf{TF-IDF baselines} already achieve ROC-AUC $>$ 0.995, indicating strong surface-level patterns in the dataset.
\item \textbf{BERT-large (3 epochs)} achieves perfect AUC = 1.0000 and is the overall best model.
\item \textbf{Ablation studies} show:
  \begin{itemize}[nosep]
  \item Learning rate: $2\times10^{-5}$ is optimal; higher rates cause slight overfitting
  \item Batch size: Smaller batches improve accuracy (BS=8: 0.9974) via implicit regularization
  \item Sequence length: 512 is marginally better than 256 (+0.0001 AUC, +0.44\% accuracy)
  \item Epochs: 3 is sufficient; 5 epochs causes degradation for both base and large models
  \end{itemize}
\item \textbf{Adversarial robustness}: The detector catches 100\% of single-pass LLM rewrites, 80\% of iterative attacks, and 96.7\% of programmatic perturbations.
\item \textbf{Error analysis}: Only 18 misclassifications out of 8,974 samples. Longer essays are classified more accurately.
\item \textbf{Scaling}: BERT-large provides perfect AUC at $3.2\times$ training cost --- diminishing returns but measurable improvement.
\end{enumerate}

\subsection{Limitations \& Future Work}
\begin{itemize}[nosep]
\item The dataset may contain artifacts that inflate performance (e.g., formatting differences)
\item More sophisticated adversarial attacks (sentence-level mixing, multiple LLM rewriting) may reduce detector confidence
\item Cross-domain generalization (different topics, different AI generators) remains untested
\end{itemize}

% ============================================================
\section*{Environment}
% ============================================================

\begin{itemize}[nosep]
\item \textbf{GPU}: NVIDIA GeForce RTX 4090 (24GB VRAM)
\item \textbf{Python}: 3.13
\item \textbf{PyTorch}: with CUDA + fp16 mixed precision
\item \textbf{Transformers}: HuggingFace
\item \textbf{Local LLM}: Ollama + DeepSeek-R1:8B
\end{itemize}

\section*{Source Code}

\begin{itemize}[nosep]
\item Notebook: \texttt{HW2.ipynb}
\item Scripts: \texttt{Baseline.py}, \texttt{BERT.py}, \texttt{LocalLLM.py}
\item GitHub: \url{https://github.com/fader2077/Recurrent-Neural-Network-and-Transformer-HW/tree/main/hw2}
\end{itemize}

\end{document}
